\centerline{\bf \Huge 9 класс}

\problem{9.1}
Разбейте какой-нибудь клетчатый квадрат на клетчатые квадратики так, чтобы не все квадратики были одинаковы, но квадратиков каждого размера было одно и то же количество.

\problem{9.2}
На доске написаны пять натуральных чисел. Оказалось, что сумма любых трёх из них делится на каждое из остальных. Обязательно ли среди этих чисел найдутся четыре равных?

\problem{9.3}
Внутри параллелограмма $A B C D$ выбрана точка $E$ так, что $A E=D E$ и $\angle A B E=90^{\circ}$. Точка $M$ - середина отрезка $B C$. Найдите угол $D M E$.

\problem{9.4}
Кондитерская фабрика выпускает $N$ сортов конфет. На Новый год фабрика подарила каждому из 1000 учеников школы подарок, содержащий по конфете нескольких сортов (составы подарков могли быть разными). Каждый ученик заметил, что для любых 11 сортов конфет он получил конфету хотя бы одного из этих сортов. Однако оказалось, что для любых двух сортов найдётся ученик, получивший конфету ровно одного из этих двух сортов. Найдите наибольшее возможное значение $N$.

\problem{9.5}
Числа $x, y$ и $z$ удовлетворяют условию $x^2+y^2+z^2=1$. Докажите, что

$$
(x-y)(y-z)(x-z) \leqslant \frac{1}{\sqrt{2}}
$$